% r03_fill_K4b: proof of (K_0), part VI (continued): the order-lambda and order-e^2 densities and the vanishing of the source-linear constant.

\begin{lemma}[order $\lambda$: the $h$-inserted figure-eight]\label{rot:lem:CK3}
$f^F_\lambda(x)=s_\lambda\,4N(N+2)\lambda_r\sum_{y\in\Lambda_F}\eta'^3G^F(x,y)^2\Gamma^F(y,y)$ with a sign $s_\lambda$ common to both flows, and $|f^X_\lambda(x')-f^Y_\lambda(\tau x')|\le C\lambda_r\eta^{1/2}$ (unit scale), i.e.\ $\le C\lambda\eta^{1/2}$ in physical units.
\end{lemma}
\begin{proof}
The order-$\lambda$ vertices are \cite[(3.18)]{K1}, $\lambda_r\sum_y\eta'^3|\phi(y)|^4$, and the order-$\lambda$ part of (3.19), $\frac12\dm_1\sum_y\eta'^3|\phi(y)|^2$ with $\dm_1=-4(N+2)\lambda_rC(0)$. The connected pairing of $|\phi(x)|^2$ with $|\phi(y)|^4$ gives $4N(N+2)G(x,y)^2G(y,y)$, with $|\phi(y)|^2$ it gives $2NG(x,y)^2$; the sum is $4N(N+2)\lambda_rG(x,y)^2[G(y,y)-C(0)]=-4N(N+2)\lambda_rG(x,y)^2\Gamma(y,y)$ (\Cref{rot:lem:Gam}(i)). Re-index the $Y$-sum by $\tau$ (Q7) and telescope: the term with $[\Gamma^X(y',y')-\Gamma^Y(\tau y',\tau y')]$ is bounded by $C_\Gamma\eta^{1/2}$ times the convergent sum $\sum_{y'}\eta'^3G^Y(\tau x',\tau y')^2\le C$; the term with $[G^X(x',y')^2-G^Y(\tau x',\tau y')^2]\Gamma^X(y',y')$ is bounded through (Q6) and (Q3): for $|u|<25\eta'$ the at most $25^3\cdot4\pi/3$ sites give $\le C\eta'$, and for $|u|\ge25\eta'$, $\sum\eta'^3C_6\eta'^{1/2}|u|^{-3/2}\cdot2C|u|^{-1}e^{-\delta'|u|}\le C\eta'^{1/2}$. Convert with $\lambda_r\delta_r^{-1}=\lambda$.
\end{proof}

\begin{lemma}[order $e^2$: kernel form]\label{rot:lem:CK4}
Let $\Pi^F[L_c,L_d](y,z)$ be the one-loop self-energy kernel with scalar loop line $L_c$ and photon line $L_d$ (exchange part from two current vertices, \cite[(3.22)]{K1} with $m=1$; seagull part from the $A^2$ vertex), $\Pi_C\defeq\Pi[C,D_C]$, $\dm_2\defeq-\hat\Pi_C(0)$ (the zero-momentum reading of \cite[(3.33)]{K1} at order $e^2$ in the free-boundary convention, the same number for both flows), and $\Pi_{C,R}(y,z)\defeq\Pi_C(y,z)+\dm_2\eta'^{-3}\delta_{yz}$, so that $\sum_z\eta'^3\Pi_{C,R}(y,z)=0$. Then, with a sign $s_e$ common to both flows,
\[
f^F_{e^2}(x)=s_e\Big\{\sum_{y,z}\eta'^6G^F(x,y)\,\Pi^F[G^F,D^F](y,z)\,G^F(z,x)+\dm_2\sum_y\eta'^3G^F(x,y)^2\Big\},
\]
$|\Pi_C(y,z)|\le Ce_r^2(|u|+\eta')^{-4}e^{-cm_r|u|}$ (Q1), and $|\Pi^X_C(y',z')-\Pi^Y_C(\tau y',\tau z')|\le Ce_r^2(\eta'/|u|)^{1/8}|u|^{-4}$ for $|u|\ge25\eta'$, $u=y-z$ (telescoping the two lines, each carrying at most two differences, with (Q2)). In each flow the exchange part equals a finite sum of terms $s_\iota\sum_{w,w'}\eta'^6L_a(x,w^+)L_b(x,w'^+)L_c(w^+,w'^+)L_d(w,w')$ with $w^+\in\{w,w+\eta'e_\mu\}$, $L_a,L_b,L_c\in\{G,\nabla_\mu G\}$, $L_d\in\{D,\nabla_\mu D\}$, exactly two differences per term and at most one per line, both along the same $\mu$ (a Euclidean inner product of two gradient vectors, \Cref{rot:lem:vert}(iii)); the $A^2$ part is frame-free.
\end{lemma}
\begin{proof}
Wick algebra at zero coarse fields; the graph with both $h$-legs at one current vertex vanishes because the current one-point function $\Im G(b_+,b_-)$ is zero. At each current vertex $\sum_s(-1)^sG(x,b_{\rm ext})G_c(b_{\rm loop},\cdot)=\eta'[(\nabla_\mu G)(x,w)G_c(w+\eta'e_\mu,\cdot)-G(x,w+\eta'e_\mu)(\nabla_\mu G_c)(w,\cdot)]$; in the (loop, loop) type the loop line carries two differences, which (Q8) moves onto $G(x,\cdot)$ or $D(\cdot,w')$ with arguments shifted by at most $\eta'$. The shifts are frame-dependent ($w+\eta'e_\mu$ against $\tau w+\eta'f_\mu$); the comparison charges one per-flow Taylor step per shifted argument, a relative gain $\eta'/\rho$ that integrates to $C\eta'(1+\log(1/\eta'))$ and is absorbed.
\end{proof}

\begin{lemma}[the $C/\Gamma$ split of the order-$e^2$ density]\label{rot:lem:CK5}
With $G=C-\Gamma$, $D=D_C-\Gamma_{\rm ph}$ and $\Pi_r\defeq\Pi[G,D]$, $f^F_{e^2}=s_e\{Z^F+A^F(x)+B^F(x)\}$ exactly, where $Z^F\defeq\sum\eta'^6C(x,y)\Pi_{C,R}(y,z)C(z,x)$, $A^F(x)\defeq\sum\eta'^6[G(x,y)G(x,z)-C(x,y)C(x,z)]\Pi_{C,R}(y,z)$ and $B^F(x)\defeq\sum\eta'^6G(x,y)G(x,z)[\Pi_r-\Pi_C](y,z)$. \textup{(i)} $Z^F$ is independent of $x$ and $Z^X=Z^Y$ (the all-$C$ part cancels exactly; it contains the logarithm of \Cref{rot:rem:log}). \textup{(ii)} $A^F(x)=\sum_{y\ne z}\eta'^6\Pi_C(y,z)[\Gamma(x,y)-\Gamma(x,z)]\{[C(x,y)-C(x,z)]-\frac12[\Gamma(x,y)-\Gamma(x,z)]\}$. \textup{(iii)} $\Pi_r-\Pi_C=-\Pi[\Gamma,D_C]-\Pi[C,\Gamma_{\rm ph}]+\Pi[\Gamma,\Gamma_{\rm ph}]-\Pi^{\rm sg}[\Gamma_{\rm ph}]$: every term has at least one internal $\Gamma$-line. The counterterm never appears alone: it is bound into $\Pi_{C,R}$, which occurs only in $Z$ and $A$. If King's scheme differs from the zero-momentum reading by a constant $\varrho=O(e_r^2)$, $f^F_{e^2}$ acquires $\varrho\sum_y\eta'^3G(x,y)^2$, whose two-flow difference is $\le C|\varrho|\eta'^{1/2}$ by the argument of \Cref{rot:lem:CK3}.
\end{lemma}
\begin{proof}
(i) is \Cref{rot:lem:ctsame}'s argument. (ii) $GG-CC=-\Gamma C-C\Gamma+\Gamma\Gamma$; symmetrise in $(y,z)$ using $\Pi(y,z)=\Pi(z,y)$ and the zero sum of $\Pi_{C,R}$ in $z$. (iii) is multilinearity.
\end{proof}

\begin{lemma}[the external-$\Gamma$ part]\label{rot:lem:CK6}
$|A^F(x)|\le Ce_r^2(1+m_r^{-1})$ and $|A^X(x')-A^Y(\tau x')|\le Ce_r^2\big[\eta^{1/4}L_\eta\log(1/\eta')+\eta'^{1/8}(1+\log(1/\eta'))^2(1+m_r^{-1})\big]$.
\end{lemma}
\begin{proof}
With $\rho=|x-y|\le\rho'=|x-z|$ and $u=y-z$: per flow $\mathsf a\defeq|\Gamma(x,y)-\Gamma(x,z)|\le C\min(|u|,1)(1+\rho)^{-1}e^{-c'm_r\rho}$ for $|u|\le\rho/2$, $\le2C(1+\rho)^{-1}e^{-c'm_r\rho}$ and $\le\sqrt3C'_\nabla|u|$ always ((Q4), $c'=c/2$); $\mathsf a_C\defeq|C(x,y)-C(x,z)|\le C|u|(\rho+\eta')^{-2}$ for $|u|\le\rho/2$, $\le2C(\rho+\eta')^{-1}$ always (Q1). Difference elements: $\delta\Pi\le Ce_r^2(\eta'/|u|)^{1/8}|u|^{-4}$ (\Cref{rot:lem:CK4}); $\delta\mathsf a\le C\min\{\eta^{1/2};\ \eta'^{1/2}(\rho^{-3/2}+\rho'^{-3/2})\ (\rho,\rho'\ge25\eta');\ A_1|u|+L_\eta|u|^2+\eta'\}$ ((Q6), \Cref{rot:cor:C}, \Cref{rot:lem:CK2}(a) with the lattice Taylor expansion of $\Gamma(x,\cdot)$ in each flow); $\delta\mathsf a_C\le C(\eta'/\rho)^{1/2}(\rho+\eta')^{-1}$ (Q2). Telescope (ii) of \Cref{rot:lem:CK5} and sum over the three regions $|u|\le\rho/2$; $|u|>\rho/2$, $\rho\le1$; $|u|>\rho/2$, $\rho>1$, bounding lattice sums by integrals and using $\sum_y\eta'^3(\rho+\eta')^{-2}e^{-cm_r\rho}\le C/m_r$, $\sum_y\eta'^3\rho^{-3}e^{-cm_r\rho}\le C\log(1/m_r)$; the regions $\rho,|u|<25\eta'$ are bounded by the per-flow sizes. The $\delta\Pi$ terms give $C\eta'^{1/8}(1+\log(1/\eta'))(1+m_r^{-1})$, the $\delta\mathsf a$ terms $C[\eta^{1/4}L_\eta\log(1/\eta')+\eta'^{1/2}\log(1/\eta')m_r^{-1/2}]$, the $\delta\mathsf a_C$ terms $C\eta'^{1/2}(1+\log(1/\eta'))^2$, and $\log(1/m_r)\le\log(1/\eta')$ absorbs the remaining logarithms.
\end{proof}

\begin{lemma}[the internal-$\Gamma$ part]\label{rot:lem:CK7}
$|B^F(x)|\le Ce_r^2$ and $|B^X(x')-B^Y(\tau x')|\le Ce_r^2\eta'^{1/6}(1+\log(1/\eta'))^2$.
\end{lemma}
\begin{proof}
By \Cref{rot:lem:CK4,rot:lem:CK5}(iii), $B^F$ is a finite sum of terms with externals $L_a,L_b\in\{G,\nabla G\}$, an internal pair with at least one $\Gamma$-type line and at most one difference per line, plus the $A^2$ terms with the photon tadpole $\Gamma_{\rm ph}(w,w)$. Per-flow sizes: $|L_a|\le C(\rho+\eta')^{-1-j_a}e^{-\delta'\rho}$ (Q3); a $C$-type internal line with $j$ differences $\le C(|u|+\eta')^{-1-j}$, a $\Gamma$-type one $\le C(1+|u|)^{-1}$ (Q1, Q4). The total scaling dimension is $6-(1+j_a)-(1+j_b)-(1+j_C)\ge1>0$ and every sub-region converges, whence the size bound; the region where $\rho$, $\rho'$ or $|u|$ is $\le M\eta'$ contributes $\le CM\eta'(1+\log(1/\eta'))$. Telescope over the lines. Outside that region the difference elements are $\delta G\le(\eta'/\rho)^{1/2}C\rho^{-1}$ (Q6), $\delta\nabla G\le(\eta'/\rho)^{1/6}\log(1/\eta)\,C\rho^{-2}$ (\Cref{rot:lem:CK2}(c)), $C$-type internal $\le(\eta'/|u|)^{1/2^{j+1}}C|u|^{-1-j}$, $j\le1$ (Q2), $\Gamma$-type internal $\delta\Gamma\le C\eta^{1/2}$, $\delta\nabla\Gamma\le A_1$ (\Cref{rot:lem:CK2}(a)). The gains $(\eta'/\rho)^\theta$ are integrated against the size integrand region by region: $\int\rho^2\rho^{-1}(\eta'/\rho)^{1/2}d\rho\le C\eta'^{1/2}$, $\int\rho^2\rho^{-2}(\eta'/\rho)^{1/6}\log(1/\eta)\,d\rho\le C\eta'^{1/6}\log(1/\eta)$, $\int|u|^2|u|^{-1-j}(\eta'/|u|)^{1/2^{j+1}}d|u|\le C\eta'^{1/4}$, and the $\Gamma$-type gain $C\eta^{1/4}L_\eta(1+|u|)$ is confined by the externals' decay; the joint sub-regions have the margin $1$ of the power counting and every gain exponent is $\le1$, so each is $\le C\eta'^{1/6}(1+\log(1/\eta'))^2$. Inside the near region the difference is at most the sum of the two per-flow sizes, $\le C\eta'(1+\log(1/\eta'))$. The photon-tadpole difference is $\le C_\Gamma\eta^{1/2}$ (Q6), and the $\eta'^2(\nabla G)^2$ term has size $\le Ce_r^2\eta'$.
\end{proof}

\begin{proposition}[the source-linear constant vanishes]\label{rot:prop:CK}
With $c_k$ as in \Cref{rot:lem:C} and $r=k-\lfloor\vr k\rfloor$, $\vr<\frac19$, for $k\ge k_0(a,m,\lmu,\vr)$
\[
|c_k|\le NC_\Gamma\,\eta^{1/2}\delta_r^{-1}+C_1\lambda\,\eta^{1/2}+C_2e^2\,\eta'^{1/8}(1+\delta_r^{-1})(1+\log(1/\eta'))^4\le C\,5^{-(1-9\vr)k/8}(1+k)^4,
\]
with constants depending on $a,m,\lmu,N,\tau$ only (not on $J$, $\psi$, $n$, the torus or $h$). Hence $c_k\to0$.
\end{proposition}
\begin{proof}
Both densities are $Y_{k-r}$-periodic (\Cref{rot:lem:C}) and the cell average of $f^Y$ equals the average of $f^Y\circ\tau$ over the matching $X$-cell (\Cref{rot:lem:vert}(i)), so $|\bar f^X-\bar f^Y|\le\sup_{x'}|f^X(x')-f^Y(\tau x')|$. By \Cref{rot:lem:S1}, $f^F=f^F_0+f^F_\lambda+f^F_{e^2}$. \emph{Term 0}: at zero coarse fields $\langle|\phi'(x)|^2\rangle=NG^F(x,x)=N(C(0)-\Gamma^F(x,x))$, so the free tadpole leaves the residue $f^F_0(x)=-N\Gamma^F(x,x)$ (the subtraction in \cite[(3.21)]{K1} removes the full-propagator expectation, not this), with matched-pair difference $\le NC_\Gamma\eta^{1/2}$ (Q6); in physical units $\Gamma$ carries $\delta_r^{-1}$, and $\eta^{1/2}\delta_r^{-1}=5^{1/2}5^{-(k-3\lfloor\vr k\rfloor)/2}\le C5^{-k/3}$. The order-$\lambda$ term is \Cref{rot:lem:CK3}, $\eta^{1/2}\le C5^{-4k/9}$; the order-$e^2$ term is \Cref{rot:lem:CK5,rot:lem:CK6,rot:lem:CK7} with $e_r^2\delta_r^{-1}=e^2$ and $m_r^{-1}=m^{-1}\delta_r^{-1}$, and $\eta'^{1/8}\delta_r^{-1}\le C5^{-(1-9\vr)k/8}$ because $\vr<\frac19$. The rate is not optimal; only $c_k\to0$ is used.
\end{proof}
